% =====================================================================
%  델파이 전문가 설문지 (독립 문서)
%  빌드: 저장소 루트에서 latexmk  →  survey/delphi-survey.pdf
% =====================================================================
\documentclass[10pt,a4paper]{article}

\usepackage{kotex}
% ---- 한글 글꼴 (나눔) + 한자 글꼴 (Noto CJK: 나눔에 없는 劍·殘·蹲踞 등) ----
\setmainhangulfont{NanumMyeongjo}[BoldFont=NanumMyeongjoBold,ItalicFont=NanumMyeongjo]
\setsanshangulfont{NanumGothic}[BoldFont=NanumGothicBold,ItalicFont=NanumGothic]
\setmonohangulfont{NanumGothicCoding}
\setmainhanjafont{Noto Serif CJK KR}[ItalicFont=Noto Serif CJK KR]
\setsanshanjafont{Noto Sans CJK KR}[ItalicFont=Noto Sans CJK KR]

\usepackage[top=22mm,bottom=22mm,left=20mm,right=20mm]{geometry}
\usepackage{amssymb,enumitem,multicol,fancyhdr,lastpage}
\setlength{\parindent}{0pt}
\setlength{\parskip}{4pt}

% =====================================================================
%  논문 정보 — 이 파일만 수정하면 표지·초록·승인서에 자동 반영됩니다.
% =====================================================================
\newcommand{\ThesisTitleKo}{논문 제목을 입력하세요}
\newcommand{\ThesisTitleEn}{Enter Your Dissertation Title Here}
\newcommand{\AuthorKo}{홍길동}
\newcommand{\AuthorEn}{Gildong Hong}
\newcommand{\AdvisorKo}{지도교수명}
\newcommand{\AdvisorEn}{Advisor Name}
\newcommand{\UniversityKo}{○○대학교 대학원}
\newcommand{\UniversityEn}{Graduate School, ○○ University}
\newcommand{\DepartmentKo}{○○학과}
\newcommand{\DepartmentEn}{Department of ○○}
\newcommand{\DegreeKo}{박사}
\newcommand{\DegreeEn}{Doctor of Philosophy}
\newcommand{\SubmitYear}{2026}
\newcommand{\SubmitMonth}{12}
\newcommand{\KeywordsKo}{키워드1, 키워드2, 키워드3}
\newcommand{\KeywordsEn}{keyword1, keyword2, keyword3}
\newcommand{\StudentID}{20XX-XXXXX}

% =====================================================================
%  델파이 설문 문항 조판 매크로
%  (설문지와 논문 부록에서 공통으로 사용 — 패키지 로드 후 \input)
% =====================================================================
\usepackage{longtable,array,colortbl,xcolor}
\usepackage[most]{tcolorbox}

\definecolor{DelphiDomain}{HTML}{1F3A5F}
\definecolor{DelphiModule}{HTML}{E8EEF5}
\definecolor{DelphiNew}{HTML}{B45309}
\definecolor{DelphiNewBg}{HTML}{FFF7ED}
\definecolor{DelphiRule}{HTML}{B8C2CC}

% ---- 설정값 (문서에서 \renewcommand / \setboolean 으로 변경) ----
\newif\ifShowProposedItems  \ShowProposedItemstrue   % [신규] 제안 문항 표시
\newif\ifShowCommentColumn  \ShowCommentColumntrue   % 문항별 의견 칸
\newif\ifShowDomainRating   \ShowDomainRatingtrue    % 영역 전체 타당성 문항
\newif\ifShowOpenBox        \ShowOpenBoxtrue         % 영역별 개방형 의견란
\providecommand{\DelphiRound}{1}                     % 조사 차수 (2 이상이면 이전 차수 결과 열 표시)

% ---- 열 너비 ----
\newlength{\DelphiNoW}      \setlength{\DelphiNoW}{10mm}
\newlength{\DelphiRateW}    \setlength{\DelphiRateW}{7.5mm}
\newlength{\DelphiCommentW} \setlength{\DelphiCommentW}{28mm}
\newlength{\DelphiPrevW}    \setlength{\DelphiPrevW}{17mm}
\newlength{\DelphiItemW}

% ---- 이전 차수 통계 (analysis/delphi_analysis.py 가 생성) ----
%   \SetPrevStat{문항번호}{중앙값}{Q1}{Q3}{CVR}
\newcommand{\SetPrevStat}[5]{%
  \expandafter\gdef\csname delphi@prev@#1\endcsname{%
    \textbf{#2}\newline{\tiny(#3--#4)}\newline{\tiny CVR #5}}}
\newcommand{\PrevStat}[1]{%
  \ifcsname delphi@prev@#1\endcsname\csname delphi@prev@#1\endcsname\else--\fi}

% ---- 내부 구현 ----
\makeatletter
\newcount\delphi@ncols
\newif\ifdelphi@prev
\newcommand{\delphi@rows}{}
\newcommand{\delphi@add}[1]{\g@addto@macro\delphi@rows{#1}}
\newcommand{\delphi@rate}{%
  & \centering ① & \centering ② & \centering ③ & \centering ④ & \centering ⑤}

\newcommand{\delphi@setup}{%
  \ifnum\DelphiRound>1 \delphi@prevtrue\else\delphi@prevfalse\fi
  \delphi@ncols=7
  \setlength{\DelphiItemW}{\dimexpr\linewidth-\DelphiNoW-5\DelphiRateW-14\tabcolsep\relax}%
  \ifShowCommentColumn
    \advance\delphi@ncols 1
    \addtolength{\DelphiItemW}{-\dimexpr\DelphiCommentW+2\tabcolsep\relax}\fi
  \ifdelphi@prev
    \advance\delphi@ncols 1
    \addtolength{\DelphiItemW}{-\dimexpr\DelphiPrevW+2\tabcolsep\relax}\fi
  \delphi@build}

% 열 정의와 머리행을 조건에 따라 미리 조립 (표 안에서 \if 를 쓰지 않기 위함)
\newcommand{\delphi@spec}{}
\newcommand{\delphi@headrow}{}
\newcommand{\delphi@hd}[1]{\color{white}\bfseries #1}
\newcommand{\delphi@build}{%
  \gdef\delphi@spec{>{\centering\arraybackslash\small}p{\DelphiNoW}
                    >{\raggedright\arraybackslash\small}p{\DelphiItemW}}%
  \gdef\delphi@headrow{\rowcolor{DelphiDomain}\delphi@hd{번호} & \delphi@hd{세부 구성요소}}%
  \ifdelphi@prev
    \g@addto@macro\delphi@spec{>{\centering\arraybackslash\footnotesize}p{\DelphiPrevW}}%
    \xdef\delphi@prevlabel{\the\numexpr\DelphiRound-1\relax}%
    \g@addto@macro\delphi@headrow{& \delphi@hd{\scriptsize \delphi@prevlabel 차 결과}\newline
                                    \color{white}\tiny Mdn (Q1--Q3)}%
  \fi
  \g@addto@macro\delphi@spec{*{5}{>{\centering\arraybackslash}p{\DelphiRateW}}}%
  \g@addto@macro\delphi@headrow{%
    & \delphi@hd{\scriptsize 전혀\newline 부적절\newline (1)}
    & \delphi@hd{\scriptsize 부적절\newline ~\newline (2)}
    & \delphi@hd{\scriptsize 보통\newline ~\newline (3)}
    & \delphi@hd{\scriptsize 적절\newline ~\newline (4)}
    & \delphi@hd{\scriptsize 매우\newline 적절\newline (5)}}%
  \ifShowCommentColumn
    \g@addto@macro\delphi@spec{>{\raggedright\arraybackslash\small}p{\DelphiCommentW}}%
    \g@addto@macro\delphi@headrow{& \delphi@hd{수정·보완 의견}}%
  \fi}

\newcommand{\delphi@head}{%
  \arrayrulecolor{DelphiDomain}\hline
  \delphi@headrow \\ \hline
  \arrayrulecolor{DelphiRule}}

% ---- 사용자 매크로 ----
\newcommand{\delphi@domnum}{}
\newcommand{\delphi@domname}{}
\newcommand{\Domain}[3]{%
  \gdef\delphi@domnum{#1}\gdef\delphi@domname{#2}%
  \gdef\delphi@rows{}%
  \delphi@setup
  \DomainHeading{#1}{#2}{#3}}

% 영역 제목 모양 (설문지/부록에서 재정의 가능)
\providecommand{\DomainHeading}[3]{%
  \par\medskip
  \begin{tcolorbox}[colback=DelphiModule,colframe=DelphiDomain,boxrule=0.6pt,
                    arc=1mm,left=2mm,right=2mm,top=1mm,bottom=1mm,
                    title={\bfseries 영역 #1.\ #2},fonttitle=\normalsize]
    \small #3
  \end{tcolorbox}}

\newcommand{\delphi@modulerow}[3]{%
  \delphi@add{\rowcolor{DelphiModule}
    \multicolumn{\the\delphi@ncols}{l}{\small\bfseries #1\ #2#3}\\ \hline}}
\newcommand{\Module}[2]{\delphi@modulerow{#1}{#2}{}}
\newcommand{\NewModule}[2]{%
  \ifShowProposedItems
    \delphi@modulerow{#1}{#2}{\ \textcolor{DelphiNew}{\scriptsize[신규 제안]}}\fi}

\newcommand{\delphi@itemrow}[4]{%
  \delphi@add{#4 #1 & \textbf{#2:} #3}%
  \ifdelphi@prev\delphi@add{& \PrevStat{#1}}\fi
  \delphi@add{\delphi@rate}%
  \ifShowCommentColumn\delphi@add{&}\fi
  \delphi@add{\\ \hline}}
\newcommand{\Item}[3]{\delphi@itemrow{#1}{#2}{#3}{}}
\newcommand{\NewItem}[3]{%
  \ifShowProposedItems
    \delphi@itemrow{#1}{#2 \textcolor{DelphiNew}{\scriptsize[신규]}}{#3}{\rowcolor{DelphiNewBg}}\fi}

\newcommand{\EndDomain}{%
  \ifShowDomainRating
    \delphi@itemrow{D\delphi@domnum}{영역 전체}%
      {위 중분류·세부 구성요소가 이 영역을 적절하게 구성하고 있다}{\rowcolor{DelphiModule}}%
  \fi
  \begingroup
  \setlength{\LTpre}{2pt}\setlength{\LTpost}{4pt}%
  \renewcommand{\arraystretch}{1.25}%
  \expandafter\longtable\expandafter{\delphi@spec}
    \delphi@head \endfirsthead
    \delphi@head \endhead
    \delphi@rows
  \endlongtable
  \endgroup
  \ifShowOpenBox\OpenCommentBox{\delphi@domnum}\fi}

% ---- 영역별 개방형 의견란 ----
\newcommand{\delphi@blank}{\par\vspace{6mm}}
\providecommand{\OpenCommentBox}[1]{%
  \begin{tcolorbox}[colback=white,colframe=DelphiRule,boxrule=0.5pt,arc=1mm,
                    left=2mm,right=2mm,top=1mm,bottom=1mm,breakable,
                    title={\small\bfseries [개방형 의견란] 영역 #1},
                    coltitle=DelphiDomain,colbacktitle=DelphiModule]
    \small
    \textbf{추가}가 필요한 구성요소:\delphi@blank
    \textbf{수정}이 필요한 구성요소 (번호 및 수정안):\delphi@blank
    \textbf{삭제}가 필요한 구성요소 (번호 및 사유):\delphi@blank
  \end{tcolorbox}}
\makeatother

% =====================================================================
%  설문지 설정 — 배포 전에 이 파일을 수정하세요.
%  (연구자·학교 정보는 ../config/metadata.tex 에서 가져옵니다)
% =====================================================================
\newcommand{\SurveyTitle}{가상현실(VR) 기반 검도훈련 콘텐츠 구성요소 개발을 위한 델파이 조사}
\renewcommand{\DelphiRound}{1}          % 조사 차수: 1, 2, 3 ...
\newcommand{\SurveyDeadline}{2026년 ○월 ○일(○)}
\newcommand{\SurveyMinutes}{약 20--25분}
\newcommand{\ResearcherPhone}{010-0000-0000}
\newcommand{\ResearcherEmail}{researcher@example.ac.kr}
\newcommand{\IRBNumber}{○○○-2026-○○○}   % 기관생명윤리위원회 승인번호 (해당 시)

% 표시 옵션
\ShowProposedItemstrue    % 연구자 추가 제안 문항([신규]) 포함 여부
\ShowCommentColumntrue    % 문항별 수정·보완 의견 칸
\ShowDomainRatingtrue     % 영역 전체 구성 타당성 문항
\ShowOpenBoxtrue          % 영역별 개방형 의견란

\IfFileExists{round\the\numexpr\DelphiRound-1\relax-stats.tex}
  {\input{round\the\numexpr\DelphiRound-1\relax-stats}}{}

\newcommand{\ChkBox}{$\square$}
\newcommand{\fillin}[1][30mm]{\underline{\hspace{#1}}}
\newcommand{\SurveySection}[1]{%
  \par\bigskip{\large\bfseries\color{DelphiDomain}#1\par}\smallskip
  {\color{DelphiDomain}\hrule height 0.8pt}\medskip}

\pagestyle{fancy}\fancyhf{}
\renewcommand{\headrulewidth}{0pt}
\fancyfoot[C]{\small\thepage\ / \pageref{LastPage}}
\fancyfoot[R]{\scriptsize\color{gray}\DelphiRound 차 델파이 조사}

\usepackage[hidelinks,unicode]{hyperref}

\begin{document}

% ------------------------------------------------------------ 인사말
\begin{center}
  {\small\color{gray}\DelphiRound 차 델파이 조사 \quad|\quad 응답자 번호: \fillin[20mm]}\par\medskip
  {\LARGE\bfseries \SurveyTitle\par}\medskip
  {\large\color{DelphiDomain}— 전문가 타당성 평가 설문지 —\par}
\end{center}
\bigskip

\begin{tcolorbox}[colback=white,colframe=DelphiDomain,boxrule=0.6pt,arc=1mm]
안녕하십니까?

바쁘신 가운데 귀한 시간을 내어 주셔서 진심으로 감사드립니다.

본 조사는 \textbf{가상현실(VR) 기반 검도훈련 콘텐츠}에 포함되어야 할 구성요소를 도출하고,
그 타당성을 검증하기 위한 \textbf{델파이(Delphi) 조사}입니다. 선행연구 분석과 전문가 면담을 통해
도출한 5개 영역의 구성요소에 대해 검도·체육교육·VR/게임 분야 전문가 여러분의 의견을 구하고자 합니다.

델파이 조사는 동일한 전문가 패널을 대상으로 \textbf{2--3회 반복}하여 의견의 합의를 이끌어내는 방법입니다.
다음 차수에서는 이번 조사의 통계 결과(중앙값, 사분위 범위 등)를 함께 제시하여 의견을 재검토하실 수 있도록
하겠습니다. 번거로우시더라도 이후 차수 조사에도 계속 참여해 주시면 감사하겠습니다.

응답 내용은 통계법 제33조에 따라 연구 목적으로만 사용되며, 개인을 식별할 수 있는 정보는 공개되지 않습니다.
응답 소요 시간은 \SurveyMinutes 이며, \textbf{\SurveyDeadline 까지} 회신해 주시기 바랍니다.

\medskip
\begin{tabular}{@{}ll}
  연구자 & \UniversityKo\ \DepartmentKo\ \DegreeKo 과정 \AuthorKo \\
  지도교수 & \AdvisorKo \\
  연락처 & \ResearcherPhone \quad / \quad \ResearcherEmail \\
  IRB 승인번호 & \IRBNumber
\end{tabular}
\end{tcolorbox}

\begin{tcolorbox}[colback=DelphiModule,colframe=DelphiRule,boxrule=0.5pt,arc=1mm]
\small\textbf{연구 참여 동의}\quad
본인은 위 연구의 목적과 내용을 이해하였으며, 자발적으로 연구에 참여하는 것에 동의합니다.\par
\hfill \ChkBox\ 동의함 \qquad \ChkBox\ 동의하지 않음 \qquad 서명: \fillin[30mm]
\end{tcolorbox}

% ------------------------------------------------------------ I. 응답자 정보
\SurveySection{I. 응답자 일반 정보}
해당하는 곳에 \checkmark\ 표시하거나 내용을 기입해 주십시오. (전문가 패널 구성의 적절성을 보고하는 데에만 사용됩니다.)

\renewcommand{\arraystretch}{1.6}
\begin{tabular}{@{}p{30mm}p{130mm}@{}}
  \textbf{1. 주 전문 분야} & \ChkBox\ 검도 지도자(사범·코치)\quad \ChkBox\ 검도 선수/심판\quad \ChkBox\ 체육교육·스포츠과학 \\
                          & \ChkBox\ VR·게임 개발/연구\quad \ChkBox\ 교육공학\quad \ChkBox\ 기타(\fillin[30mm]) \\
  \textbf{2. 검도 단위} & \ChkBox\ 없음\quad \ChkBox\ 1--3단\quad \ChkBox\ 4--5단\quad \ChkBox\ 6단 이상 (\fillin[10mm]단) \\
  \textbf{3. 경력} & 검도 수련 \fillin[12mm]년 \quad 지도 \fillin[12mm]년 \quad 해당 분야 연구·개발 \fillin[12mm]년 \\
  \textbf{4. VR 경험} & \ChkBox\ 없음\quad \ChkBox\ 체험해 본 적 있음\quad \ChkBox\ 정기적으로 사용\quad \ChkBox\ VR 콘텐츠 개발·연구 경험 \\
  \textbf{5. 최종 학력} & \ChkBox\ 학사\quad \ChkBox\ 석사\quad \ChkBox\ 박사수료\quad \ChkBox\ 박사 \\
  \textbf{6. 소속/직위} & \fillin[80mm]
\end{tabular}
\renewcommand{\arraystretch}{1}

% ------------------------------------------------------------ II. 작성 방법
\SurveySection{II. 설문 작성 방법}
아래 제시된 각 구성요소가 \textbf{“가상현실(VR) 기반 검도훈련 콘텐츠에 포함되어야 할 필요성 및 적절성이
어느 정도인가”}에 대해 5점 척도로 평가하여 해당 번호에 \textbf{○ 또는 \checkmark} 표시해 주십시오.

\begin{center}
\renewcommand{\arraystretch}{1.3}
\begin{tabular}{|c|c|c|c|c|}
  \hline\rowcolor{DelphiModule}
  1점 & 2점 & 3점 & 4점 & 5점 \\ \hline
  \small 전혀 적절하지 않음 & \small 적절하지 않음 & \small 보통 & \small 적절함 & \small 매우 적절함 \\
  \small (불필요) & & & & \small (필수) \\ \hline
\end{tabular}
\end{center}

\begin{itemize}[leftmargin=5mm,itemsep=1pt]
  \item 각 영역의 마지막 \colorbox{DelphiModule}{\small 영역 전체} 문항은 해당 영역의 중분류·세부 구성요소 체계 전체가 적절한지를 평가합니다.
  \item 표현이 모호하거나 수정이 필요한 문항은 \textbf{수정·보완 의견} 칸에 간단히 적어 주십시오.
  \item 도출된 항목 외에 추가/수정/삭제가 필요하다고 판단되는 구성요소는 각 영역 하단의
        \textbf{[개방형 의견란]}에 자유롭게 작성해 주십시오.
\ifShowProposedItems
  \item \colorbox{DelphiNewBg}{\textcolor{DelphiNew}{\small[신규]}} 표시 문항은 검도 경기·심판 규정 및 VR 안전 지침을
        검토하여 연구자가 추가로 제안한 구성요소입니다. 다른 문항과 동일하게 평가해 주십시오.
\fi
\ifnum\DelphiRound>1
  \item \textbf{\the\numexpr\DelphiRound-1\relax 차 결과} 열에는 이전 차수 전체 응답의 중앙값(Mdn)과 사분위 범위(Q1--Q3),
        내용타당도 비율(CVR)이 제시되어 있습니다. 이를 참고하여 의견을 재검토해 주시고,
        본인의 응답이 사분위 범위를 벗어나는 경우 수정·보완 의견 칸에 그 이유를 적어 주십시오.
\fi
\end{itemize}

% ------------------------------------------------------------ III. 평가
\newpage
\SurveySection{III. VR 검도훈련 콘텐츠 구성요소 타당성 평가}
% =====================================================================
%  VR 검도훈련 콘텐츠 구성요소 — 델파이 설문 문항 (단일 원본)
%  설문지(survey/delphi-survey.tex)와 논문 부록이 모두 이 파일을 사용합니다.
%
%   \Domain{번호}{영역명}{영역 설명}      대분류
%   \Module{번호}{중분류명}                중분류 (표 안의 구분 행)
%   \Item{번호}{구성요소명}{설명}          세부 구성요소 (평정 문항)
%   \NewItem{번호}{구성요소명}{설명}       연구자 추가 제안 문항 ([신규] 표시,
%                                          \ShowProposedItems 가 false 면 숨김)
%   \EndDomain                            영역 종료 (영역 타당성 + 개방형 의견란)
% =====================================================================

% ---------------------------------------------------------------------
\Domain{1}{하드웨어 및 인터랙션 (Hardware \& Interaction)}
  {가상 공간에서 실제 검도의 신체 감각, 칼의 물리적 궤적 및 반동을 구현하기 위한 기술적·장비적 요소}

\Module{1.1}{모션 트래킹}
\Item{1.1.1}{9축/6DoF 죽도 트래킹}{죽도의 위치, 속도, 칼날 각도를 실시간 추적}
\Item{1.1.2}{HMD 시선 추적}{사용자의 시선 방향 및 주시점 인식}
\Item{1.1.3}{하체 발놀림(Foot) 추적}{스텝, 발밟기 등 하체 동작 트래킹}

\Module{1.2}{햅틱 피드백}
\Item{1.2.1}{타격 순간 반동 피드백}{격자 부위 타격 시 죽도 장비에 물리적 진동/충격 전달}
\Item{1.2.2}{교검(交劍) 저항 피드백}{상대 죽도와 칼이 맞닿거나 쳐누를 때 물리적 저항감 제공}

\Module{1.3}{공간 음향}
\Item{1.3.1}{3D 입체 음향}{타격음, 상대 기합 소리, 경기장/도장 울림 등 공간 음향 구현}
\EndDomain

% ---------------------------------------------------------------------
\Domain{2}{검도 무도론 및 교수학습 (Kendo Pedagogy)}
  {검도 고유의 수련 체계, 거리감(간합), 무도적 가치를 VR 환경에 맞춰 구조화한 콘텐츠 요소}

\Module{2.1}{수련 단계 모드}
\Item{2.1.1}{기본 수련}{중단자세, 기본 발놀림, 공간치기 반복 수련 모드}
\Item{2.1.2}{응용/격권 수련}{머리·손목·허리 연속 치기 및 받아치기/쳐내기 기법 모드}
\Item{2.1.3}{자유 대련}{AI 또는 멀티플레이어 기반 대적 시뮬레이션 모드}

\Module{2.2}{간합 및 AI}
\Item{2.2.1}{일간의 거리(一間之間) 시각화}{유효 타격이 가능한 시공간적 거리 안내 표시}
\Item{2.2.2}{대적 AI 엔진}{유저의 거리 조절에 맞춰 전진/후진 및 응대 반격을 수행하는 AI}

\Module{2.3}{기세 및 기합}
\Item{2.3.1}{기합 음성 인지}{마이크 입력을 통해 타격 순간의 기합 소리 크기 및 타이밍 인식}

\NewModule{2.4}{예법 및 무도 가치}
\NewItem{2.4.1}{예법(禮法) 수련}{입장·례(禮)·준거(蹲踞)·납도 등 검도 예법 절차 안내 및 수행 여부 확인}
\EndDomain

% ---------------------------------------------------------------------
\Domain{3}{동작 평가 및 피드백 (Assessment \& Feedback)}
  {수련자의 동작이 검도의 정격(正擊) 규칙 및 원리에 부합하는지 측정·평가하는 요소}

\Module{3.1}{심검체 일치도}
\Item{3.1.1}{심검체(心劍體) 동기화 측정}{시선(心)-죽도 궤적(劍)-발밟기(體)의 타임스탬프 일치 판정}

\Module{3.2}{정격 판정}
\Item{3.2.1}{칼날 각도(刃筋) 검증}{타격 시 죽도 날지름이 유효 범위(45도 이내)로 들어왔는지 판정}
\Item{3.2.2}{유효 격자부위 판정}{머리, 손목, 허리, 찌르기 부위에 정확히 타격되었는지 측정}
\NewItem{3.2.3}{타돌부(物打) 판정}{죽도의 유효 타돌부(선단 약 1/4 부위)로 타격했는지 판정}
\NewItem{3.2.4}{잔심(殘心) 판정}{타격 후 자세·기세·거리를 유지하며 다음 동작에 대비했는지 판정}

\Module{3.3}{다차원 피드백}
\Item{3.3.1}{3D 슬로우 모션 리플레이}{타격 순간을 360도 및 속도 조절로 재확인하는 기능}
\Item{3.3.2}{수련 데이터 시각화}{타격 속도(m/s), 반응 시간, 죽도 궤적 흔들림 다이어그램 제공}
\EndDomain

% ---------------------------------------------------------------------
\Domain{4}{시스템 및 환경 (UI/UX \& System)}
  {훈련자의 몰입감 유지, 개인 맞춤화 및 연속적인 학습 관리를 지원하는 보조 요소}

\Module{4.1}{수련 환경 구축}
\Item{4.1.1}{가상 훈련 공간 선택}{기본 도장, 체육관, 전국 대회 경기장 등 가상 환경 제공}

\Module{4.2}{아바타 커스텀}
\Item{4.2.1}{신체 치수 보정}{사용자의 키, 팔 길이, 죽도 규격에 맞춘 1:1 아바타 보정}

\Module{4.3}{학습 관리(LMS)}
\Item{4.3.1}{수련 리포트 제공}{누적 수련 시간, 부위별 타격 성공률, 심검체 일치도 추이 그래프 제공}
\NewItem{4.3.2}{지도자 모니터링 및 원격 코칭}{지도자가 수련자의 기록을 열람하고 코멘트·과제를 제공하는 기능}

\NewModule{4.4}{안전 및 사용 편의}
\NewItem{4.4.1}{안전 경계 및 충돌 방지}{실제 공간의 안전 경계(가디언) 설정 및 죽도 휘두르기 공간 확보 안내}
\NewItem{4.4.2}{사이버멀미 및 피로 관리}{프레임률 유지, 이동 방식 선택, 휴식 알림 등 VR 멀미·신체 피로 관리}
\EndDomain

% ---------------------------------------------------------------------
\Domain{5}{게임 설계 및 몰입 (Game Design \& Engagement)}
  {수련자의 지속적인 훈련 동기 유발, 몰입(Flow) 유지 및 학습 효과 극대화를 위한 게임화 요소}

\Module{5.1}{동기부여 및 보상}
\Item{5.1.1}{단/급 승단 체계}{수련 성과 기반 단/급 수여 및 진급 칭호 제공}
\Item{5.1.2}{수련 미션/퀘스트}{일일 수련 목표 설정 및 미션 달성 시 보상 부여}

\Module{5.2}{난이도 및 밸런싱}
\Item{5.2.1}{동적 난이도 조절(DDA)}{수련자 숙련도 및 정격률에 맞춘 AI 공격 속도/패턴 실시간 조절}
\Item{5.2.2}{수련 커브 디자인}{입문부터 고단자 수련까지 난이도 상승 곡선(Learning Curve) 단계화}

\Module{5.3}{몰입 및 연출}
\Item{5.3.1}{정격 타격 이펙트}{무도적 품위를 유지하는 수준의 정격 시각/청각 피드백 연출}

\Module{5.4}{경쟁 및 비동기 수련}
\Item{5.4.1}{랭킹 및 리더보드}{정격 정확도 및 심검체 일치도 점수 기반 랭킹 시스템}
\Item{5.4.2}{고단자 고스트(Ghost) 대전}{고단자의 실제 모션 데이터를 기반으로 한 비동기 대적 수련}
\EndDomain


% ------------------------------------------------------------ IV. 상대적 중요도
\SurveySection{IV. 영역별 상대적 중요도}
VR 검도훈련 콘텐츠 개발에 있어 각 영역이 차지하는 상대적 중요도를 \textbf{합계가 100점이 되도록} 배분해 주십시오.
(콘텐츠 개발 우선순위 및 영역별 가중치 산출에 활용됩니다.)

\begin{center}
\renewcommand{\arraystretch}{1.6}
\begin{tabular}{|p{95mm}|c|}
  \hline\rowcolor{DelphiModule}
  \textbf{영역} & \textbf{배분 점수} \\ \hline
  1. 하드웨어 및 인터랙션 & \fillin[15mm] 점 \\ \hline
  2. 검도 무도론 및 교수학습 & \fillin[15mm] 점 \\ \hline
  3. 동작 평가 및 피드백 & \fillin[15mm] 점 \\ \hline
  4. 시스템 및 환경 & \fillin[15mm] 점 \\ \hline
  5. 게임 설계 및 몰입 & \fillin[15mm] 점 \\ \hline
  \rowcolor{DelphiModule}\multicolumn{1}{|r|}{\textbf{합계}} & \textbf{100 점} \\ \hline
\end{tabular}
\end{center}

% ------------------------------------------------------------ V. 종합 의견
\SurveySection{V. 종합 의견}
\begin{tcolorbox}[colback=white,colframe=DelphiRule,boxrule=0.5pt,arc=1mm]
\small
1. 영역(대분류) 구성 자체에 대해 추가·통합·분리가 필요하다고 생각하시는 부분이 있다면 적어 주십시오.\par\vspace{16mm}
2. VR 검도훈련 콘텐츠가 실제 도장 수련을 보완하기 위해 가장 중요하다고 생각하시는 점은 무엇입니까?\par\vspace{16mm}
3. 기타 본 연구에 대한 제언이 있으시면 자유롭게 적어 주십시오.\par\vspace{16mm}
\end{tcolorbox}

\begin{center}
\bigskip
{\large 끝까지 응답해 주셔서 진심으로 감사드립니다.}
\end{center}

\end{document}
